\scanpage{210}

Si en particulier pour tt $S'$, $X_{S'}^{p/S'}\deleted{\unsure{\not\sim|\not\ni}}\xleftarrow{\sim} X_{S'}^{p/S'\,!}$, ce qui est le cas en particulier si $X/S$ est \uline{lisse}, on trouve un isom

\[
P^{cip}_{X/S}\simeq \underline{\mathrm{Pic}}^{\deleted{\unsure{cp}}}_{X^{p/S}/S}\simeq\underline{\mathrm{Pic}}_{X/S}^{(p/S)}\tag*{(3.4)}
\]

L'homomorphisme canonique $P^{cip}_{X/S}\to\underline{\mathrm{Pic}}_{X/S}$ n'est alors autre que la Verschiebung

\[
V\colon\underline{\mathrm{Pic}}_{X/S}^{p/S}\longrightarrow\underline{\mathrm{Pic}}_{X/S}\tag*{(3.5)}
\]

[\unsure{encore|avouer|ainsi} qu'il faudrait définir $V$ avec une généralité suffisante pour que ce qu'on vient de dire ait tjrs un sens \ldots]. On trouve donc, en \unsure{vertu} \deleted{\unsure{en}} (3.2), un isom.

\[
\mathrm{Ker}(V_{\underline{\mathrm{Pic}}_{X/S}})\simeq
\left\{\begin{aligned}
&\text{faisceau des \deleted{formes}\inserted{sections} $\omega$ de $f_*^{\mathrm{fl}}(\Omega^1_{X/S})$ telles}\\
&\text{que $d\omega=0$, $\widetilde C\omega=\widetilde\omega$, modulo le sous-}\\
&\text{faisceau des $df/f$, $f$ ds $f_*^{\mathrm{fl}}(\underline O^*_{X/S})$}
\end{aligned}\right.\tag*{\unsure{(3.6)|(2.6)}}
\]

Quand on inverse sur $(\mathrm{Sch})/S$ les \deleted{\unsure{immersions}} morphismes entiers surjectifs radiciels, \inserted{[en fait, il suffit \unsure{pour} certains très particuliers]}, de \deleted{\unsure{sorte}} que $\underline{\mathrm{Pic}}_{X/S}\xrightarrow F\underline{\mathrm{Pic}}_{X/S}^{p/S}$ devient inversible, on trouve que le \deleted{\unsure{deuxième}} premier mbr devient $\,{}_p\underline{\mathrm{Pic}}_{X/S}$ \ldots.

\uline{Comparaison} \deleted{des} $P^{cip\,\unsure{0|o}}_{X/S}$ \uline{et} $P^{cip\,0}_{\mathrm{Alb}^0_{X/S}}$ Plaçons nous pour simplifier dans le cas $X/S$ lisse, de sorte que (3.4) est valable, et supposons que $\underline{\mathrm{Pic}}^0_{X/S}$ soit un schéma abélien $B$. Comme on aura $(\underline{\mathrm{Pic}}^{\deleted{\unsure{p/S}}}_{X/S})^{(p/S)\,0}=(\underline{\mathrm{Pic}}^0_{X/S})^{p/S}=B^{(p/S)}$, \deleted{$\simeq\underline{\mathrm{Pic}}^0_{\unsure{X}}$} \deleted{\unsure{on trouve}} et qu'on a aussi $B\simeq\underline{\mathrm{Pic}}^0_{A^1/S}=\underline{\mathrm{Pic}}^0_{A/S}$,

% archive sheet number: 209
